% Generated by scripts/tikz_reviews.py using the compiled book's labels.
\pdfbookmark[0]{Chapter 18}{comparison-chapter-18}
\ReviewFigure{fig-convex-examples}
{Convexity of functions and sets}
{figures/convexity/convex-examples.png}
{figures/tikz/convex-analysis-convex-examples.pdf}
{The epigraph definition of convexity and the distinction between convexity and strict convexity.}
{The same original also appears in the smooth-optimization chapter. This reconstruction is shared in design with that chapter: segments exhibit failure of convexity, an affine or flat boundary segment, and strict convexity. The epigraph is shaded above each function. The functions and sets are schematic examples rather than named formulas from the source.}
{convex-analysis-convex-examples}
{Complete figure}
{18.1}
\ReviewFigure{fig-review-convex-analysis-subdifferential}
{Supporting affine functions at a kink}
{figures/convexity/subdifferential.png}
{figures/tikz/convex-analysis-subdifferential.pdf}
{Definition of the subdifferential as all supporting slopes.}
{The original fan of supporting lines is reconstructed at a nonsmooth point of a convex function. Every drawn slope belongs to the interval between the one-sided derivatives. The diagram labels the supporting expression with a separate evaluation variable z, avoiding confusion with the base point x.}
{convex-analysis-subdifferential}
{Complete figure}
{18.2}
\ReviewFigure{fig-subdiff-l1}
{Graphs of two subdifferentials}
{figures/convexity/subdiff-l1.png}
{figures/tikz/convex-analysis-subdiff-l1.pdf}
{The subdifferential of absolute value and a piecewise affine convex function.}
{Interpretation to check: At each kink the entire closed interval of intermediate slopes is drawn, not just its endpoints. The second function uses representative slopes m1<m2<0<m3 at knots a<0 and 0; the original does not specify numerical slopes. The lower panels are set-valued graphs, so their vertical segments are essential.}
{convex-analysis-subdiff-l1}
{Complete figure}
{18.3}
\ReviewFigure{fig-review-convex-analysis-normal-cones}
{Normal cones at boundary, interior, and exterior points}
{figures/convexity/normal-cone.png}
{figures/tikz/convex-analysis-normal-cones.pdf}
{The normal cone is the subdifferential of the indicator of a convex set.}
{The convex body is rendered schematically as a polygon so its normal directions are exact. Edge interiors have outward normal rays, corners have cones, interior points have only the zero normal, and exterior points have an empty normal cone. The arrows depict directions and are translated to their base points; they do not bound the cone lengths.}
{convex-analysis-normal-cones}
{Complete figure}
{18.4}
